% kuleuventheme2 by Janez Kren, September 2017, janez.kren@kuleuven.be, based on:
% kuleuventheme 1.3 by Roland Pastorino, 2013 roland.pastorino@kuleuven.be / www.rolandpastorino.com

\PassOptionsToPackage{usenames, table, dvipsnames}{xcolor} % 
\documentclass[11pt,t]{beamer}
\usetheme{kuleuven2}	%THEME OPTIONS for LOGO: kul (default), kulak, lrd,    ; OPTIONS for TITLE PAGE: normal (default), sedes


%%% OTHER SETTINGS
\usefonttheme[onlymath]{serif}			% math font with serifs, delete to make it sans-serif
\setbeamertemplate{footline}[body] 		% delete this line to remove footline bar on all frames
%\usepackage[orientation=landscape,size=custom,width=16,height=9,scale=0.5,debug]{beamerposter} %enable for widescreen 16:9 ratio
%\titlegraphic{ \includegraphics[width=.2\paperwidth]{mytitlepagepic.png} } %optional title page image


%%% ADDED BY ALEX (start)
\usepackage[english]{babel}
\usepackage{amsfonts}
\usepackage{amssymb}

\usepackage{pythonhighlight}

\usepackage[orientation=landscape,size=custom, width=16,height=9,scale=0.5,debug]{beamerposter} 

\definecolor{mygray}{rgb}{0.8,0.8,0.8}
\newcommand{\inline}[1]{ \colorbox{mygray}{\lstinline{#1}}}

\usepackage{xcolor}

\usepackage{layouts} % To get text/page width

\usepackage{listings}
\definecolor{codegreen}{rgb}{0,0.6,0}
\definecolor{codegray}{rgb}{0.5,0.5,0.5}
\definecolor{codepurple}{rgb}{0.58,0,0.82}
% \definecolor{backcolour}{rgb}{0.95,0.95,0.92}
\definecolor{backcolour}{rgb}{0.93,0.93,0.95}
\lstdefinestyle{mystyle}{
    backgroundcolor=\color{backcolour},   
    commentstyle=\color{codegreen},
    keywordstyle=\color{magenta},
    numberstyle=\tiny\color{codegray},
    stringstyle=\color{codepurple},
    basicstyle=\ttfamily\footnotesize,
    breakatwhitespace=false,         
    breaklines=true,                 
    captionpos=b,                    
    keepspaces=true,                 
    numbers=left,                    
    numbersep=5pt,                  
    showspaces=false,                
    showstringspaces=false,
    showtabs=false,                  
    tabsize=2
}
\lstset{style=mystyle}

%%% ADDED BY ALEX (stop)


%%% TITLE PAGE INFO:
\title[Geavanceerde Computerarchitectuur Project Introduction]{Geavanceerde Computerarchitectuur Project Introduction} %[]] will appear in footline
\subtitle{Goals, examples, and grading}

\author{Thomas Feys, {\footnotesize thomas.feys@kueluven.be}}
\institute{DRAMCO -- WaveCore -- ESAT}
\date{Academic year 2024/25}




\begin{document}
\csname beamer@calculateheadfoot\endcsname %recalculate head and foot dimension


 %%
 %%  0. TITLE PAGE and TABLE OF CONTENT
 %%
% Title page
\begin{frame}[plain,noframenumbering]
	\titlepage
\end{frame}
	

% Table of Contents
\begin{frame}{Outline}
	\hfill	{\large \parbox{.961\textwidth}{\tableofcontents[hideothersubsections]}}
\end{frame}







%%
%%  SECTION 1 - INFO
%%
\section{Recap from the first session}

\begin{frame}{Organization of the lab and docs} 
    \begin{figure}[t]
        \centering
        \begin{center}               
            \resizebox{0.98\columnwidth}{!}{
\begin{tikzpicture}[]

\node[align=center, anchor=center] (gca) at (0, 0) {GCA lab};



\node[align=center, anchor=center] (docs) at (0, 1.5) {Docs\\(Toledo)};
\node[align=center, anchor=center, font=\footnotesize, text=black!60!white] at (-2, 2.25) {Prog. guide};
\node[align=center, anchor=center, font=\footnotesize, text=black!60!white] at ( 0, 2.5) {Cheat sheet};
\node[align=center, anchor=center, font=\footnotesize, text=black!60!white] at ( 2, 2.25) {tasks.json};



\node[align=center, anchor=center] (sess) at (3, -1) {5 Sessions};
\node[align=center, anchor=center, font=\footnotesize, text=black!60!white] at ( 3, 0 ) {Kernel};
\node[align=center, anchor=center, font=\footnotesize, text=black!60!white] at ( 4.5, 0 ) {Race c.};
\node[align=center, anchor=center, font=\footnotesize, text=black!60!white] at ( 5, -0.75 ) {Warps};
\node[align=center, anchor=center, font=\footnotesize, text=black!60!white] at ( 5, -1.5 ) {Memory};
\node[align=center, anchor=center, font=\footnotesize, text=black!60!white] at ( 4.5, -2.25 ) {Dynamic p.};
\node[align=center, anchor=center, font=\footnotesize, text=black!60!white] at ( 3, -2.25 ) {Async.};



\node[align=center, anchor=center] (proj) at (-3, -1) {Project};
% \node[align=center, anchor=center, font=\footnotesize, text=black!60!white] at ( -3, 0 ) {Graphics};
\node[align=center, anchor=center, font=\footnotesize, text=black!60!white] at ( -4.5, 0 ) {\dots};
\node[align=center, anchor=center, font=\footnotesize, text=black!60!white] at ( -5, -0.75 ) {Neural Networks};
\node[align=center, anchor=center, font=\footnotesize, text=black!60!white] at ( -5, -1.5 ) {Encryption};
\node[align=center, anchor=center, font=\footnotesize, text=black!60!white] at ( -4.5, -2.25 ) {DSP};
\node[align=center, anchor=center, font=\footnotesize, text=black!60!white] at ( -3, -2.25 ) {Graphics};


\draw[black] (gca.south west) rectangle (gca.north east);

\draw[] (gca.south west) -- (proj.north east);
\draw[] (gca.south east) -- (sess.north west);
\draw[] (gca.north) -- (docs.south);


\end{tikzpicture}
}
        \end{center}
    \end{figure}
\end{frame}

\begin{frame}{Responsibilities and evaluation} 
    \begin{figure}[t]
        \centering
        \begin{center}               
            \resizebox{0.98\columnwidth}{!}{
\begin{tikzpicture}[]

\foreach \i in {0,...,1} {
    \node[anchor=center, align=center, font=\footnotesize] (lr\i) at (-5.5, -1.2*\i) {Short report};
    \draw[black] (lr\i.south west) rectangle (lr\i.north east);
    
    \node[anchor=center, align=center, font=\footnotesize] (rr\i) at (-2, -1.2*\i) {Short report};
    \draw[black] (rr\i.south west) rectangle (rr\i.north east);
}
\node[anchor=center, align=center, font=\footnotesize] (rr2) at (-2, -1.2*2) {Short report};
\draw[black] (rr2.south west) rectangle (rr2.north east);

\node[anchor=center, align=center, font=\footnotesize] (pr) at (3.5,0) {Project report};
\draw[black] (pr.south west) rectangle (pr.north east);
    
\node[anchor=center, align=center, font=\footnotesize] (pp) at (3.5,-1.2) {Project Presentation};
\draw[black] (pp.south west) rectangle (pp.north east);
    
\node[anchor=center, align=center, font=\footnotesize] (cr) at (3.5,-2.4) {Code release};
\draw[black] (cr.south west) rectangle (cr.north east);


\draw[dashed] (0, 1.5) -- (0, -3);

\node[] at (-3.5, 1) {Lab sesions, 50 \%};
\node[] at ( 3.5, 1) {Project, 50 \%};


\end{tikzpicture}
}
        \end{center}
    \end{figure}
\end{frame}


\section{Project}


\begin{frame}{Goal and organization} 

\begin{itemize}\setlength\itemsep{0.5em}
\item The goal of the project is to tackle a computational challenge using the GPU.
\item This should be something that aligns with your interests -- to support your thesis or another lab; to implement something for a hobby / private project; to boost your portfolio and get a better job.
\item We will allocate 5 lab sessions to development and 1 session to the presentations.
\end{itemize}

\end{frame}


\begin{frame}{Implementation} 

\begin{itemize}\setlength\itemsep{0.75em}
\item Use the concepts we have gotten acquainted with so far and feel free to explore other paradigms. You can use another programming language and also build on another library -- keep in mind that you should still understand and describe the intricacies of the implementation.
\item At the lab, you can use the B230 PCs, an Nvidia Jetson Nano ("Raspberry with a GPU"), or the two Nvidia Teslas in the server. You are free to use other devices if you can/want.
\end{itemize}

\end{frame}


\begin{frame}{Some topic examples} 
In case you don't know where to start, here are a few biased examples:
\vspace{2mm}
\begin{itemize}\setlength\itemsep{0.5em}
    \item An RF \textbf{ray tracer} to determine how Wi-Fi waves propagate in (and outside of) your kot.
    \item An \textbf{optimization algorithm} to determine where to put the Wi-Fi router for optimal coverage (considering either only yourself or all inhabitants).
    \item \textbf{Data compression} on an edge device that collects sensor samples and analyzes/compresses/transforms these in order to send the information over a low-throughput link.
    \item Real-time \textbf{statistical analysis} of a video stream.
\end{itemize}
\end{frame}

\begin{frame}{Some topic examples} 
    In case you don't know where to start, here are a few biased examples:
    \vspace{2mm}
    \begin{itemize}\setlength\itemsep{0.5em}
        \item Implement a neural network using CUDA: fully connected layers, convolution layers, pooling layers, ... (training can be done in python)
        \item Implement game of life on CUDA (cellular automaton)
        \item Anomaly Detection in Time Series Data
        \item Implement Clustering algorithms (K-means, DBSCAN, ...)
        \item Implement and optimize a support vector machine using CUDA
        \item Implement backpropagation for simple MLP
    \end{itemize}
    \end{frame}



\begin{frame}{Grading -- guidelines based on thesis evaluation matrix} 

\vspace{-2mm}

Implementation
{\small
\begin{itemize}
    \item The implementation works, while clearly addressing the tackled problem
    \item Includes the concepts introduced in the lab or sufficiently motivates other approaches
\end{itemize}
}

\vspace{1mm}

Report (no min/max length -- try to be concise)
{\small
\begin{itemize}
    \item Clearly describes the problem and motivates the usage of the GPU
    \item Introduces the implementation and methods in detail (includes references to the repository or code listings)
    \item Highlights the shortcomings and possible improvements
\end{itemize}
}

\vspace{1mm}

Presentation (up to 10 min + 10 min for Q\&A)
{\small
\begin{itemize}
    \item Clearly describes the problem and motivates the usage of the GPU
    \item Provides a high-level overview of the implementation
    \item Highlights the shortcomings and possible improvements
\end{itemize}
}

\end{frame}



\begin{frame}{Important dates}

Please submit the reports on time, so we can close off the lab on the day of the presentations (focus on the thesis / exams afterwards).

\vspace{6mm}

TO BE DETERMINED

The important dates for {\color{olive}Monday's} and {\color{Plum}Friday's} group are:
\vspace{2mm}
\begin{itemize}\setlength\itemsep{0.75em}
    \item Project sessions -- {\color{olive}17/03, 24/03, 28/04, 05/05, 12/05}
    \item Hand in the reports -- {\color{olive}14/05}
    \item Presentations and finish -- {\color{olive}19/05} 
\end{itemize}

\vspace{4mm}

{\footnotesize the session of 31/03 will be cancelled!.}

\end{frame}


\end{document}
